\section*{Lab 5 (Decision Trees) Implementation Tasks}
\textbf{Lab 5 focus:} supervised rule learning, purity and impurity, split search, recursive tree growth, interpretability, and overfitting control.

\medskip
\textbf{Implementation policy (strict):}
\begin{itemize}
  \item Implement the decision tree core algorithm with NumPy only.
  \item Do not use a library decision tree implementation for the main task.
  \item Plotting libraries are allowed for visual inspection.
  \item External libraries are allowed only for data loading if explicitly needed.
  \item Work through the notebook in order.
\end{itemize}

\section*{Core Equations And Rules}
\begin{itemize}
  \item Supervised dataset:
  \[
    D=\{(x_i,y_i)\}_{i=1}^{n},
    \qquad
    x_i\in\mathbb{R}^d,
    \qquad
    y_i\in\{1,\dots,C\}.
  \]
  \item Class proportion at node $m$:
  \[
    p_{mc}
    =
    \frac{1}{n_m}
    \sum_{(x_i,y_i)\in D_m}\mathbf{1}\{y_i=c\}.
  \]
  \item Gini impurity:
  \[
    I_G(D_m)=1-\sum_{c=1}^{C}p_{mc}^{2}.
  \]
  \item Entropy impurity:
  \[
    I_H(D_m)=-\sum_{c=1}^{C}p_{mc}\log_2 p_{mc}.
  \]
  \item Pure node: $I_G(D_m)=0$; larger values indicate mixed classes.
  \item Leaf prediction (majority class):
  \[
    \hat{y}=\arg\max_{c\in\{1,\dots,C\}} p_{mc}.
  \]
  \item Candidate split:
  \[
    D_L=\{(x_i,y_i)\in D_m:x_{ij}\leq t\},
    \qquad
    D_R=\{(x_i,y_i)\in D_m:x_{ij}>t\}.
  \]
  \item Impurity decrease:
  \[
    \Delta I(j,t)
    =
    I(D_m)
    -
    \frac{|D_L|}{|D_m|}I(D_L)
    -
    \frac{|D_R|}{|D_m|}I(D_R).
  \]
\end{itemize}

\section*{Data Files And Preparation}
\begin{itemize}
  \item \texttt{decision\_trees.ipynb}: guided notebook with the assignment steps.
  \item \texttt{advanced\_ai/trees/decision\_tree.py}: starter implementation.
  \item \texttt{prepare\_datasets.py}: generates the prepared data bundles.
  \item \texttt{scripts/prepare\_all\_data.sh}: shell entrypoint for dataset preparation.
  \item \texttt{scripts/check\_prepare\_all\_data.sh}: verification script for the prepared bundles.
  \item \texttt{toy\_tree.py}, \texttt{iris\_tree.py}, \texttt{depth\_and\_overfitting.py}: task runners.
  \item \texttt{collect\_submission.ipynb} and \texttt{collectSubmission.sh}: submission helpers.
\end{itemize}

\medskip
Expected prepared files:
\begin{itemize}
  \item \texttt{datasets/toy\_tree.npz}
  \item \texttt{datasets/iris\_tree.npz}
  \item \texttt{datasets/depth\_study.npz}
\end{itemize}

\medskip
The scripts use the conda environment \texttt{ai\_courses}.

\section*{Roadmap}
\begin{itemize}
  \item Start with Step 1: prepare and verify the datasets.
  \item After that, continue through the steps in order.
  \item Each step introduces only the task you need next.
\end{itemize}

\section*{Step 1: Prepare the Datasets}
\begin{itemize}
  \item From \texttt{assignment5\_decision\_trees/}, run:
  \[
    \texttt{bash scripts/prepare\_all\_data.sh}
  \]
  \item Then run:
  \[
    \texttt{bash scripts/check\_prepare\_all\_data.sh}
  \]
  \item Confirm that all three expected \texttt{.npz} bundles exist.
\end{itemize}

\section*{Step 2: Warm-up by Hand}
\begin{itemize}
  \item Use a parent node with class counts $(6,4)$.
  \item Compute the class proportions.
  \item Compute the parent Gini impurity.
  \item Evaluate the candidate split with child counts $(5,0)$ and $(1,4)$.
  \item Report the weighted child impurity and the impurity decrease.
\end{itemize}

\section*{Step 3: Implement the Decision Tree Core}
\begin{enumerate}[label=\textbf{T\arabic*.}]
  \item \textbf{Gini impurity}
\begin{itemize}
  \item Open \texttt{advanced\_ai/trees/decision\_tree.py}.
  \item Fill in \texttt{\_gini}.
  \item Verify the pure case gives $I_G=0$ and the balanced binary case gives $I_G=0.5$.
\end{itemize}

  \item \textbf{Split scoring}
\begin{itemize}
  \item Fill in \texttt{\_split\_gain}.
  \item Use the weighted impurity decrease $\Delta I(j,t)$ from the lecture.
  \item Respect \texttt{min\_samples\_leaf}: a split is invalid if either child is too small.
  \item Invalid splits should have no positive gain.
\end{itemize}

  \item \textbf{Candidate split search}
\begin{itemize}
  \item Fill in \texttt{\_best\_split}.
  \item Search all features.
  \item For each feature, test thresholds halfway between consecutive unique feature values.
  \item Return the feature index, threshold, and impurity decrease of the best valid split.
\end{itemize}

  \item \textbf{Recursive tree growth}
\begin{itemize}
  \item Fill in \texttt{\_build\_node}.
  \item Stop when the node is pure, maximum depth is reached, the node is too small, or no split improves impurity enough.
  \item Each leaf should store the majority-class prediction of the examples that reach it.
\end{itemize}

  \item \textbf{Prediction and path explanation}
\begin{itemize}
  \item Use the provided \texttt{predict}, \texttt{depth}, \texttt{n\_leaves}, and \texttt{explain\_one\_path} methods after \texttt{\_build\_node} is complete.
  \item Do not change their public behavior unless you document the reason.
\end{itemize}
\end{enumerate}

\section*{Step 4: Run the Self-Checks}
\begin{itemize}
  \item After implementing the core class, run:
  \[
    \texttt{python -m advanced\_ai.self\_checks}
  \]
  \item Continue only after every check passes.
\end{itemize}

\section*{Step 5: Toy Dataset}
\begin{itemize}
  \item Run \texttt{toy\_tree.py}.
  \item Report the tree depth, number of leaves, and training accuracy.
  \item Save the decision-region plot.
  \item Explain one root-to-leaf path in plain language.
  \item Identify the first split and state which feature and threshold it uses.
\end{itemize}

\section*{Step 6: Iris Data}
\begin{itemize}
  \item Run \texttt{iris\_tree.py}.
  \item Compare raw and standardized features.
  \item Sweep over maximum depth.
  \item Record training accuracy, validation accuracy, and the number of leaves.
  \item Explain whether standardization changed the tree behavior and why this differs from k-NN, k-means, and DBSCAN.
\end{itemize}

\section*{Step 7: Depth and Overfitting}
\begin{itemize}
  \item Run \texttt{depth\_and\_overfitting.py}.
  \item Compare shallow and deep decision regions.
  \item Explain why training accuracy can improve while validation accuracy stops improving.
  \item Choose one depth as your final model and justify the choice using validation accuracy and interpretability.
\end{itemize}

\section*{Step 8: Optional Extension}
\begin{itemize}
  \item Try entropy instead of Gini impurity.
  \item Try minimum leaf sizes from \texttt{MIN\_LEAF\_CHOICES}.
  \item Report whether the selected root split changes.
\end{itemize}

\section*{Submission}
\begin{itemize}
  \item Submit the notebook, the completed starter files, and the generated output files from the required runners.
  \item Use \texttt{collect\_submission.ipynb} or \texttt{collectSubmission.sh} to bundle the submission.
\end{itemize}
